% GENERATED FILE. Edit canonical lessons, then run npm run generate:books.
% canonical-source-hash: fnv1a64:681ca915a18c487b
% canonical-lessons: ZH-C30-zenmeyang, ZH-C30-hen, ZH-W30-hen, ZH-C30-henhao, ZH-C30-mang, ZH-C30-lei, ZH-C30-haixing, ZH-C30-practice

\chapter{How Are You?}
\label{ch:zh-how-are-you}
\hlchaptermodality{30}

\begin{chapteropening}
\textbf{By the end of this chapter:} \emph{I can ask how somebody is and answer that I am fine, busy, tired or so-so.}

The everyday way to ask how somebody is, and four answers to it, built on one new written character: \zh{很}, the plain link Mandarin puts in front of an adjective.
\end{chapteropening}

\section[zěnmeyàng]{zěnmeyàng — how are you doing?}

\label{lesson:ZH-C30-zenmeyang}



\begin{quote}
A new chapter, and it opens with a question you can ask anybody you know.

\begin{itemize}
\raggedright
  \item \emph{Recall:} say \emph{nǐ ne} — \textbf{R1}, the line that turned the passage one lesson back
  \item \emph{Recall:} say \emph{nǐ xǐhuan shénme} — \textbf{R2}, the first real retrieval, five to fifteen lessons back
  \item \emph{Recall:} say \emph{wǒ yǒu yí ge míngzi} — \textbf{R3}, durable, twenty to sixty lessons back
  \item \emph{Recall:} say \emph{zǎoshang hǎo} — \textbf{R4}, recognition at distance, eighty lessons back or more
\end{itemize}
\end{quote}

\begin{sounds}
\emph{zěnmeyàng} — \textbf{a dip, a light syllable, then a fall.}

\emph{zěn} is third tone. \emph{me} is the same toneless syllable that ends \zh{什么}. \emph{yàng} is fourth tone, and it carries the weight of the word.

\begin{quote}
\emph{nǐ zěnmeyàng} — How are you doing?
\end{quote}

Say it once slowly and once at speed: \textbf{zěn me yàng}.
\end{sounds}

\begin{culture}
You already have \textbf{\zh{你好吗}} \emph{nǐ hǎo ma}, and it is correct. It is the question textbooks open with.

Between people who know each other, the one you will hear far more often is \emph{nǐ zěnmeyàng}: \textquotedblleft{}how are things with you?\textquotedblright{} On its own, \emph{zěnmeyàng} asks how something is, and putting \zh{你} in front makes the something you.

So keep both. \zh{你好吗} is the careful one; \emph{nǐ zěnmeyàng} is the one a friend asks.
\end{culture}

\begin{cousinweb}
Mandarin builds \emph{zěnmeyàng} from three syllables. \emph{me} is the light \zh{么} from \zh{什么}, a character you know. The characters for \emph{zěn} and \emph{yàng} come later; this chapter keeps them to the ear.

No dictionary we use records an older form of \emph{zěn} or \emph{yàng}, so the history of the word is not settled, and where \zh{么} comes from is not settled either.
\end{cousinweb}

\subsection*{Your turn}
Say these aloud:

\begin{itemize}
\raggedright
  \item \emph{zěnmeyàng} — dip, light, fall
  \item \emph{nǐ zěnmeyàng}
  \item \zh{你好吗}, then the everyday one, \emph{nǐ zěnmeyàng}
\end{itemize}

\begin{tcolorbox}[breakable,colback=teal!4,colframe=teal!35!black,title={Before you move on}]
What tones does \emph{zěnmeyàng} carry? (\textbf{Third, none, fourth.}) Which question would a friend ask you? (\emph{\textbf{nǐ zěnmeyàng}} — \zh{你好吗} is the one textbooks open with.)
\end{tcolorbox}
\section[\texorpdfstring{\zh{很} (hěn)}{很 (hěn)}]{\zh{很} — the word in front of an adjective}

\label{lesson:ZH-C30-hen}



\begin{quote}
The question is in place. This lesson gives you the word every answer needs.

\begin{itemize}
\raggedright
  \item \emph{Recall:} say \emph{nǐ zěnmeyàng} — \textbf{R1}, one lesson back, before it decays
  \item \emph{Recall:} say \emph{xǐhuan} — \textbf{R2}, the first real retrieval, five to fifteen lessons back
  \item \emph{Recall:} say \emph{wǒ de shū} — \textbf{R3}, durable, twenty to sixty lessons back
  \item \emph{Recall:} say \emph{hǎo} — \textbf{R4}, recognition at distance, the first word in this book
\end{itemize}
\end{quote}

\begin{sounds}
\textbf{\zh{很}} is \emph{hěn} — \textbf{third tone}, the dip, the same as \zh{好}.

Its pieces, in words: on the left, the person piece from \zh{你} with one more left-falling stroke above it; on the right, a piece read \emph{gěn}. The next lesson puts it in your hand.
\end{sounds}

\begin{grammarlens}[title={Grammar Lens: no \zh{是}, and a plain \zh{很} instead}]
\begin{quote}
\textbf{\zh{我很好}} \emph{wǒ hěn hǎo} — I'm fine.
\end{quote}

In Mandarin an adjective is the whole predicate: there is no \zh{是} in \textquotedblleft{}I am fine\textquotedblright{}, and a plain \zh{很} stands in front of \zh{好} instead. Here \zh{很} has lost most of its \emph{very}, so \zh{我很好} is \textquotedblleft{}I'm fine\textquotedblright{}, while a bare \zh{我好} sounds like the start of a contrast — I am fine, but somebody else is not.
\end{grammarlens}

\begin{cousinweb}
\textbf{\zh{很}} is an old character: the Guangyun records it as \emph{honX} (in Baxter's transcription), glossed \textquotedblleft{}defiant\textquotedblright{}. How that word came to mean \textquotedblleft{}very\textquotedblright{} is not settled.
\end{cousinweb}

\subsection*{Your turn}
Say these aloud:

\begin{itemize}
\raggedright
  \item \emph{hěn} — dip it
  \item \zh{我很好}
  \item \emph{nǐ zěnmeyàng} — \zh{我很好}
\end{itemize}

\begin{tcolorbox}[breakable,colback=teal!4,colframe=teal!35!black,title={Before you move on}]
Does \zh{是} go in front of an adjective? (\textbf{No — a plain \zh{很} does.}) What does \zh{我很好} mean? (\textbf{I'm fine} — \zh{很} has lost most of its \emph{very}.)
\end{tcolorbox}
\section[\texorpdfstring{\zh{很} (hěn)}{很 (hěn)}]{\zh{很} — nine strokes, two pieces}

\label{lesson:ZH-W30-hen}



\begin{quote}
Three characters first, read cold, before the new piece arrives.

\begin{itemize}
\raggedright
  \item \emph{Recall:} say \emph{hěn} — \textbf{R1}, one lesson back, before it decays
  \item \emph{Recall:} say \emph{shíbā suì} — \textbf{R2}, the first real retrieval, five to fifteen lessons back
  \item \emph{Recall:} say \emph{yǒu} — \textbf{R3}, durable, twenty to sixty lessons back
\end{itemize}

\begin{quote}
\zh{人} \zh{亻} \zh{你}
\end{quote}

\emph{Read it:} each one aloud, with nothing said first — \emph{rén}, the person piece, \emph{nǐ}
\end{quote}

\subsection*{Script — two pieces, left before right}
\hlblockfigure{\detokenize{figures/ZH-W30-hen-filmstrip.pdf}}{How \zh{很} is written, stroke by stroke}

\begin{quote}
\zh{很}
\end{quote}

\textbf{Nine strokes, but two pieces.} Do not count to nine. Read the character across, left piece first, the way you write \zh{你}.

\textbf{The left piece, \zh{彳}, three strokes.} Its lower two are the person piece \zh{亻} you have written since \zh{你}: a left-falling stroke, then a vertical down to the foot. \zh{彳} adds one more short left-falling stroke above them. So: a short fall, a longer fall, then the vertical.

\textbf{The right piece, \zh{艮}, six strokes.} A horizontal that turns down the right side of a box without lifting. A short horizontal inside it. A horizontal that closes the box. Then the left side, down to the foot, flicking up to the right. A short left-falling stroke on the right, and last the long right-falling stroke.

3 + 6 = nine, and the first three are a shape your hand already knows.

\subsection*{Your turn}
\begin{itemize}
\raggedright
  \item \emph{Write it:} \zh{亻}, then \zh{彳} — the same two strokes with one more fall above them
  \item \emph{Write it:} \zh{很}, all nine, with the model in front of you
  \item \emph{Cover:} write it again from memory, and count the pieces, not the strokes
\end{itemize}

\begin{tcolorbox}[breakable,colback=teal!4,colframe=teal!35!black,title={Before you move on}]
How many strokes does \zh{很} have? (\textbf{Nine}.) How many pieces? (\textbf{Two} — \zh{彳}, then \zh{艮}.) Which piece is the person piece with one stroke added? (\textbf{\zh{彳}, on the left}.)

Source: \href{https://raw.githubusercontent.com/chanind/hanzi-writer-data/68d10a4b21150cae5e1ebbd223eed289cf32d90c/data/\%E5\%BE\%88.json}{Hanzi Writer Data \zh{很}, pinned stroke paths and medians}.
\end{tcolorbox}
\section[\texorpdfstring{\zh{很好} (hěn hǎo)}{很好 (hěn hǎo)}]{\zh{很好} — I'm fine}

\label{lesson:ZH-C30-henhao}



\begin{quote}
Nine strokes in your hand. Now they become an answer.

\begin{itemize}
\raggedright
  \item \emph{Recall:} say \emph{hěn} — \textbf{R1}, the character one lesson back, before it decays
  \item \emph{Recall:} say \emph{duì} — \textbf{R2}, the first real retrieval, five to fifteen lessons back
  \item \emph{Recall:} say \emph{nǐ hǎo} — \textbf{R4}, recognition at distance, the first word in this book
\end{itemize}
\end{quote}

\begin{sounds}
\textbf{\zh{很好}} is \emph{hěn hǎo} — \textbf{two third tones side by side.}

You met this in the first chapter with \zh{你好}, and the same rule fires again: the first of the two rises. Spoken, it is \emph{hén hǎo}. Written, it stays \emph{hěn hǎo}, and the characters do not change at all.

Say it slowly once, then at speed, and let \zh{很} rise into \zh{好}.
\end{sounds}

\begin{culture}
\begin{quote}
\textbf{\zh{我很好}} \emph{wǒ hěn hǎo} — I'm fine.
\end{quote}

That is the answer to \emph{nǐ zěnmeyàng} and to \zh{你好吗}. It means \textquotedblleft{}I'm fine\textquotedblright{}, not \textquotedblleft{}I'm very well\textquotedblright{}: \zh{很} is there because \zh{好} is an adjective and wants it in front. Put weight on \zh{很}, and the \emph{very} comes back.
\end{culture}

\begin{tcolorbox}[breakable,skin=enhanced,colback=orange!6,colframe=orange!45!black,boxrule=0.5pt,arc=1mm,left=6pt,right=6pt,top=4pt,bottom=4pt,fonttitle=\bfseries,title={Read this}]
\begin{quote}
\zh{我很好你呢}
\end{quote}

Five characters, and nothing on the page marks where one word stops. Find the words before reading on.

\zh{我} — \zh{很好} — \zh{你} — \zh{呢}. Four words; \zh{很好} is the only one with two characters.
\end{tcolorbox}

\begin{cousinweb}
\textbf{\zh{很好}} is built from \zh{很} and \zh{好}. \zh{好} is inherited from Middle Chinese \emph{xawX} (Guangyun, in Baxter's transcription); \zh{很} is the character the Guangyun records as \emph{honX}, \textquotedblleft{}defiant\textquotedblright{}, and how it came to mean \textquotedblleft{}very\textquotedblright{} is not settled.
\end{cousinweb}

\subsection*{Your turn}
\begin{itemize}
\raggedright
  \item \emph{Say it:} \emph{hén hǎo} — let the first one rise
  \item \emph{Say it:} \emph{nǐ zěnmeyàng} — \zh{我很好}
  \item \emph{Write it:} \zh{很好}
\end{itemize}

\begin{tcolorbox}[breakable,colback=teal!4,colframe=teal!35!black,title={Before you move on}]
Which of the two syllables in \zh{很好} rises when you say it? (\textbf{The first, \zh{很}.})
\end{tcolorbox}
\section[máng]{máng — busy}

\label{lesson:ZH-C30-mang}



\begin{quote}
One answer so far. This lesson adds a second, and it is not \textquotedblleft{}fine\textquotedblright{}.

\begin{itemize}
\raggedright
  \item \emph{Recall:} say \emph{wǒ hěn hǎo} — \textbf{R1}, one lesson back, before it decays
  \item \emph{Recall:} say \emph{nǐ zěnmeyàng} — \textbf{R2}, the question that opened this chapter
  \item \emph{Recall:} say \emph{yǒu} — \textbf{R3}, durable, twenty to sixty lessons back
  \item \emph{Recall:} say \emph{zǎoshang} — \textbf{R4}, recognition at distance, eighty lessons back or more
\end{itemize}
\end{quote}

\begin{sounds}
\emph{máng} — \textbf{second tone}, rising, like a surprised \textquotedblleft{}huh?\textquotedblright{}. Its character comes later; this chapter keeps it to the ear.
\end{sounds}

\begin{grammarlens}[title={Grammar Lens: the same shape as \zh{我很好}}]
\emph{máng} is an adjective, so it works exactly like \zh{好}: no \zh{是}, a plain \zh{很} in front.

\begin{quote}
\emph{wǒ hěn máng} — I'm busy. \emph{zǎoshang wǒ hěn máng} — I'm busy in the mornings.
\end{quote}

\emph{zǎoshang} goes first, before the person: the time is what the sentence is about, then comes what happens in it.
\end{grammarlens}

\begin{cousinweb}
\emph{máng} is probably inherited from Middle Chinese \emph{mang} (Guangyun, in Baxter's transcription), where it is glossed \textquotedblleft{}afraid\textquotedblright{}. Flustered becoming busy is a likely path for the sense, but it is not settled.
\end{cousinweb}

\subsection*{Your turn}
Say these aloud:

\begin{itemize}
\raggedright
  \item \emph{máng} — let it rise
  \item \emph{wǒ hěn máng}
  \item \emph{zǎoshang wǒ hěn máng}
\end{itemize}

\begin{tcolorbox}[breakable,colback=teal!4,colframe=teal!35!black,title={Before you move on}]
What tone is \emph{máng}? (\textbf{Second, rising.}) What goes in front of it? (\textbf{A plain \zh{很}}, as in front of \zh{好}.) Where does \emph{zǎoshang} go? (\textbf{First, before the person.})
\end{tcolorbox}
\section[lèi]{lèi — tired}

\label{lesson:ZH-C30-lei}



\begin{quote}
Busy all morning, and now you are tired. You step into a lift, and the panel reads:

\begin{quote}
\zh{九}
\end{quote}

\emph{Say it:} \emph{jiǔ} — the ninth floor, read off the wall, the way you first read floor numbers

\begin{itemize}
\raggedright
  \item \emph{Recall:} say \emph{máng} — \textbf{R1}, one lesson back, before it decays
  \item \emph{Recall:} say \emph{xǐhuan} — \textbf{R2}, the first real retrieval, five to fifteen lessons back
  \item \emph{Recall:} say \emph{yí ge} — \textbf{R3}, durable, twenty to sixty lessons back
\end{itemize}
\end{quote}

\begin{sounds}
\emph{lèi} — \textbf{fourth tone}, a short clean fall. Its character comes later; this chapter keeps it to the ear.
\end{sounds}

\begin{grammarlens}[title={Grammar Lens: a statement, then \zh{吗}}]
\emph{lèi} is an adjective too. Ask with \zh{吗} on the end, and answer with \zh{很}:

\begin{quote}
\emph{nǐ lèi ma} — Are you tired? \emph{wǒ hěn lèi} — I'm tired.
\end{quote}

The question is the statement with \zh{吗} added; nothing moves.
\end{grammarlens}

\begin{cousinweb}
The character for \emph{lèi} is in the Guangyun, read \emph{ljweX} and \emph{ljweH} (Baxter's transcription), but glossed \textquotedblleft{}accumulate\textquotedblright{} and \textquotedblleft{}implicate\textquotedblright{}. \textquotedblleft{}Tired\textquotedblright{} is not among its senses there, so how the word came to mean it is not settled.
\end{cousinweb}

\subsection*{Your turn}
Say these aloud:

\begin{itemize}
\raggedright
  \item \emph{lèi} — fall
  \item \emph{nǐ lèi ma}
  \item \emph{wǒ hěn lèi}
\end{itemize}

\begin{tcolorbox}[breakable,colback=teal!4,colframe=teal!35!black,title={Before you move on}]
What tone is \emph{lèi}? (\textbf{Fourth, falling.}) How do you ask \textquotedblleft{}are you tired?\textquotedblright{} (\emph{\textbf{nǐ lèi ma}} — the statement, then \zh{吗}.) And the lift panel? (\textbf{\zh{九}, \emph{jiǔ}.})
\end{tcolorbox}
\section[hái xíng]{hái xíng — not bad}

\label{lesson:ZH-C30-haixing}



\begin{quote}
Fine, busy, tired. Most days are none of those, and this lesson is for them.

\begin{itemize}
\raggedright
  \item \emph{Recall:} say \emph{lèi} — \textbf{R1}, one lesson back, before it decays
  \item \emph{Recall:} say \emph{duì} — \textbf{R2}, the first real retrieval, five to fifteen lessons back
  \item \emph{Recall:} say \emph{nǐ hé wǒ} — \textbf{R3}, durable, twenty to sixty lessons back
  \item \emph{Recall:} say \emph{bù} — \textbf{R4}, recognition at distance, eighty lessons back or more
\end{itemize}
\end{quote}

\begin{tcolorbox}[breakable,skin=enhanced,colback=orange!6,colframe=orange!45!black,boxrule=0.5pt,arc=1mm,left=6pt,right=6pt,top=4pt,bottom=4pt,fonttitle=\bfseries,title={Read this}]
\begin{quote}
\zh{早上好} \zh{你好吗} \zh{我很好} \zh{你呢}
\end{quote}

\emph{Read it:} the three lines aloud as one exchange, without stopping between them

The last line hands the question back. This lesson gives you a third way to answer it.
\end{tcolorbox}

\begin{sounds}
\emph{hái xíng} — \textbf{two rising tones.}

\emph{hái} is \textquotedblleft{}still\textquotedblright{}; \emph{xíng} is \textquotedblleft{}OK, it works\textquotedblright{}. Together: \textquotedblleft{}still OK\textquotedblright{} — not bad. Its characters come later; this chapter keeps the word to the ear.

\begin{quote}
\emph{hái xíng, nǐ ne} — Not bad. And you?
\end{quote}
\end{sounds}

\begin{culture}
\zh{很好} says things are good. \emph{hái xíng} says they are all right, and it is the answer people give when they are not going to complain. Said with a shrug, it can mean \textquotedblleft{}could be better\textquotedblright{}.
\end{culture}

\begin{cousinweb}
\emph{hái xíng} is built from \emph{hái} \textquotedblleft{}still\textquotedblright{} and \emph{xíng} \textquotedblleft{}all right\textquotedblright{}. Both characters are in the Guangyun (Baxter's transcription): the one for \emph{hái} as \emph{hwaen}, \textquotedblleft{}return\textquotedblright{}, and the one for \emph{xíng} as \emph{haeng}, \textquotedblleft{}walk, go\textquotedblright{}. How they came to mean \textquotedblleft{}still\textquotedblright{} and \textquotedblleft{}all right\textquotedblright{} is not settled.
\end{cousinweb}

\subsection*{Your turn}
Say these aloud:

\begin{itemize}
\raggedright
  \item \emph{hái xíng} — rise, rise
  \item \emph{hái xíng, nǐ ne}
\end{itemize}

\begin{tcolorbox}[breakable,colback=teal!4,colframe=teal!35!black,title={Before you move on}]
What tones does \emph{hái xíng} carry? (\textbf{Two rising tones.}) Which answer is between \emph{hěn hǎo} and \emph{hěn lèi}? (\emph{\textbf{hái xíng}} — not bad.)
\end{tcolorbox}
\section[Practice]{nǐ zěnmeyàng — wǒ hěn hǎo}

\label{lesson:ZH-C30-practice}



\begin{quote}
The chapter closes here, with a whole exchange.

\begin{itemize}
\raggedright
  \item \emph{Recall:} say \emph{hái xíng} — \textbf{R1}, one lesson back, before it decays
  \item \emph{Recall:} say \emph{wǒ xǐhuan kàn shū} — \textbf{R2}, the first real retrieval, five to fifteen lessons back
  \item \emph{Recall:} say \emph{wǒ shì zhōngguórén} — \textbf{R3}, durable, twenty to sixty lessons back
  \item \emph{Recall:} say \emph{zǎoshang hǎo} — \textbf{R4}, recognition at distance, eighty lessons back or more
\end{itemize}
\end{quote}

\subsection*{The exchange}
\begin{quote}
\textbf{A:} \zh{早上好}
\end{quote}

\begin{quote}
\emph{zǎoshang hǎo} — good morning.
\end{quote}

\begin{quote}
\textbf{B:} \zh{早上好} — \emph{nǐ zěnmeyàng}
\end{quote}

\begin{quote}
\emph{zǎoshang hǎo, nǐ zěnmeyàng} — good morning. How are you doing?
\end{quote}

\begin{quote}
\textbf{A:} \zh{我很好} \zh{你呢}
\end{quote}

\begin{quote}
\emph{wǒ hěn hǎo, nǐ ne} — I'm fine. And you?
\end{quote}

\begin{quote}
\textbf{B:} \emph{hái xíng, wǒ hěn máng} — not bad. I'm busy.
\end{quote}

\begin{quote}
\textbf{A:} \emph{nǐ lèi ma} — are you tired?
\end{quote}

\begin{quote}
\textbf{B:} \emph{hěn lèi, zàijiàn} — I am. Goodbye.
\end{quote}

\begin{quote}
\textbf{A:} \zh{再见}
\end{quote}

\begin{quote}
\emph{zàijiàn} — goodbye.
\end{quote}

The words whose characters this book has not written yet stay in pinyin; you know them by ear, and here the ear carries them. The morning greeting is one of the first things this book gave you, and this is the first exchange to carry on past it.

\begin{tcolorbox}[breakable,skin=enhanced,colback=orange!6,colframe=orange!45!black,boxrule=0.5pt,arc=1mm,left=6pt,right=6pt,top=4pt,bottom=4pt,fonttitle=\bfseries,title={Read this}]
\emph{Read it:} the exchange through once without stopping, then answer the question below
\end{tcolorbox}

\begin{grammarlens}[title={What you've built — four answers to one question}]
One question, \emph{nǐ zěnmeyàng}, and four answers, every one of them with a plain \zh{很} or none at all:

\begin{itemize}
\raggedright
  \item \emph{wǒ hěn hǎo} — I'm fine.
  \item \emph{hái xíng} — not bad.
  \item \emph{wǒ hěn máng} — I'm busy.
  \item \emph{wǒ hěn lèi} — I'm tired.
\end{itemize}

\textbf{Listen.} \emph{Nǐ zěnmeyàng?}

\textbf{Speak.} Answer it four ways, then hand it back with \zh{你呢}.

\emph{Write it:} \zh{很} — \zh{彳} first, then \zh{艮}
\end{grammarlens}

\subsection*{Your turn}
\begin{itemize}
\raggedright
  \item \emph{Say it:} \emph{nǐ zěnmeyàng} — then answer with \emph{hái xíng}
  \item \emph{Say it:} the whole exchange, both parts
  \item \emph{Write it:} \zh{我很好}
\end{itemize}

\begin{tcolorbox}[breakable,colback=teal!4,colframe=teal!35!black,title={Before you move on}]
Ask a friend how they are doing. (\emph{\textbf{nǐ zěnmeyàng}}.) Say you are fine. (\textbf{\zh{我很好}} — and \zh{很} rises into \zh{好}.) What goes in front of an adjective? (\textbf{A plain \zh{很}}, and no \zh{是}.)
\end{tcolorbox}
